To assess the robustness of the bottleneck selection and the classical reference methods, all models were evaluated with the same stratified grouped 5-fold cross-validation protocol, using machining run as the grouping variable. The autoencoder showed stable performance across bottleneck dimensions: the best mean F1 was obtained with a 128-dimensional bottleneck and the global MAE reconstruction score (precision 0.798, recall 0.964, F1 0.868, PR-AUC 0.849), while the 64-dimensional bottleneck with global MSE achieved a nearly identical F1 score (precision 0.811, recall 0.933, F1 0.865, PR-AUC 0.855). In terms of threshold-free PR-AUC, the 16-dimensional bottleneck with global MAE was marginally strongest among the autoencoder variants (PR-AUC 0.857), but the differences between the 16-, 32-, and 64-dimensional configurations were small relative to the observed cross-validation variability. The descriptor-based baselines were competitive under the same fold assignments, with isolation forest achieving the highest mean F1 and PR-AUC (precision 0.873, recall 0.924, F1 0.891, PR-AUC 0.897), followed by one-class SVM (precision 0.883, recall 0.878, F1 0.880, PR-AUC 0.887). However, these differences to the best autoencoder configurations were modest and within the cross-validation variability intervals, so they should not be interpreted as clear statistical superiority. PCA reconstruction produced very high recall (0.982), but at substantially lower precision (0.625), resulting in a lower F1 score (0.749) and PR-AUC (0.637). Overall, the grouped-CV results indicate that the proposed autoencoder is robust with respect to the exact bottleneck dimensionality, while the classical descriptor-based methods provide strong but not clearly superior reference baselines on this small turning dataset.
